\documentclass[11pt]{article}
\usepackage[UTF8]{ctex}
\usepackage{amsmath,amssymb,amsthm}
\usepackage{etoolbox}

% 1. 自定义定理样式实现标题后强制换行
\newtheoremstyle{break}% 样式名称
  {\topsep}{\topsep}% 环境上下间距
  {\normalfont}% 正文字体
  {}% 缩进量
  {\bfseries}{.}% 头部字体和标点
  {\newline}% 头部后的间距：设置为 \newline 即可实现完美的换行效果
  {}% 头部时间/附加参数规范

% 应用新样式并声明常用的定理类环境
\theoremstyle{break}
\newtheorem{theorem}{定理}[section]
\newtheorem{lemma}[theorem]{引理}
\newtheorem{corollary}[theorem]{推论}
\newtheorem{definition}[theorem]{定义}

% 2. 利用 etoolbox 宏包动态修补 proof 环境，使其标题后自动换行
\patchcmd{\proof}{\item[\hskip\labelsep\itshape#1\@addpunct{.}]}{\item[\hskip\labelsep\itshape#1\@addpunct{.}]\hfill\par}{}{}

\begin{document}
\section{多项式压缩与固定折叠和集大小}
\subsection{摘要}
\textbf{定义}（$N(h,k)$ 的极小覆盖性质）对于固定的正整数 $h$，令 $N(h,k)$ 为满足如下条件的最小整数 $N$：任意由 $k$ 个整数组成的集合所达到的基数 $|hA|$，均能由 $\{1,\dots,N\}$ 的某个 $k$ 元子集实现。

\textbf{定理}（$N(h,k)$ 的多项式上界）对于任意固定的 $h$，均有
\[ N(h,k) \le k^{C_h}, \]
其中 $C_h$ 为仅依赖于 $h$ 的正常数。

\textbf{注记}（证明架构与新工具）本结果建立在 Rajagopal 关于 $|hA|$ 可能取值范围的固定 $h$ 定理之上。本文的核心创新在于，在 Rajagopal 的大值构造框架内，采用多项式直径替换了原有的稀疏几何块。该替换操作严格在低次 Hilbert 函数层面执行；值得指出的是，当 $h$ 固定时，低次 Hilbert 函数构成了刻画 $h$ 重和集的唯一相关数据。
\section{引言}

对于有限集 $A \subset \mathbb{Z}$ 及正整数 $h$，记
\[
hA = \{a_1 + \cdots + a_h : a_i \in A\}
\]
（允许元素重复选取）。
\textbf{定义 1.1}（$h$ 重和） 上述集合 $hA$ 称为 $A$ 的 $h$ 重和。

\textbf{定义 1.2}（值域集 $R(h,k)$） 定义
\[
R(h,k) = \{|hA| : A \subset \mathbb{Z}, |A| = k\}.
\]

\textbf{定义 1.3}（最小表示数 $N(h,k)$） 令 $N(h,k)$ 为满足以下条件的最小正整数 $N$：
\[
R(h,k) = \{|hA| : A \subseteq \{1,\dots,N\}, |A| = k\}.
\]

\textbf{注记 1.4} 等价地，研究可在区间 $[0, N-1]$ 上进行，最终结果通过平移即可得到。

本文旨在证明下述多项式压缩界。
\textbf{定理 1.1.} 对于任意固定的 $h \ge 1$，存在常数 $C$，使得对所有 $k \ge 2$，有
\[ N(h,k) \le k^{C h}. \]
\textbf{证明.} 该证明基于 Rajagopal 关于固定折叠和集大小取值范围的定理。令
\[ M_{h,k} = h^{k-h+1}, \quad K_{h,k} = \binom{h+k-1}{h}, \]
并定义
\[ \Delta_{h,k} = \prod_{\ell=0}^{\min(h,k)-3} [M_{h,k} + \ell h + 1, M_{h,k} + \ell h + (h-2-\ell)]. \]
此处及下文，$[u,v]$ 表示满足 $u \le n \le v$ 的整数 $n$ 构成的集合；若 $u > v$，则该区间为空集。
\end{document}